\chapter{Conclusions} \label{chap:conclusiones}

In this work we have simulated the 2D Ising model on a square lattice with periodic boundary conditions, comparing three Monte Carlo algorithms (Metropolis, Glauber and Wolff) for $L\in\{16,32,64,128\}$. Our aims were twofold: to verify that all three reproduce equilibrium physics and to compare their dynamical behaviour. In equilibrium the three algorithms agree, as required when sampling the same Boltzmann distribution. The Binder-cumulant crossings yield an estimate of $T_c$ compatible with Onsager's result ($U_4^*\approx0.61$), and the FSS collapse of $\chi$ and $\xi_L/L$ using critical exponents confirms that the data belong to the universality class of the 2D Ising model. The main caveat arises near $T_c$ for observables sensitive to autocorrelation: Metropolis and Glauber display critical slowing down, so a clean convergence of $T_c$ with $L$ is observed only for Wolff, and the peak position of the collapsed susceptibility shifts to $x\approx1$ for local algorithms while Wolff's peak sits near $x\approx1.5$, which we regard as the more reliable value.

The essential difference between the algorithms lies in dynamics as measured by the integrated autocorrelation time $\tau_{\mathrm{int}}$. Metropolis and Glauber suffer critical slowing down: the peak $\tau_{\mathrm{int}}$ near $T_c$ grows with $L$ consistent with $z\approx2$, and Glauber is somewhat slower than Metropolis (factor \~1.5–2) because Metropolis typically has higher acceptance rates and thus decorrelates faster. Wolff effectively removes critical slowing down: its $\tau_{\mathrm{int}}$ is of order unity across the range (small $z$), so obtaining independent samples near $T_c$ costs orders of magnitude less than with local updates — a decisive advantage for larger lattices.

As a natural extension, it would be interesting to repeat the comparison for the three-dimensional Ising model, where no exact solution exists and critical exponents are known only numerically, now that the simulation and data-analysis pipeline has been validated against the exact 2D results.

\newpage
\clearpage
\begin{otherlanguage}{english}
\noindent\textbf{\large Conclusions (English)}
\vspace{0.5em}

In this work we have simulated the 2D Ising model on a square lattice with periodic boundaries, comparing three Monte Carlo algorithms (Metropolis, Glauber and Wolff) for $L\in\{16,32,64,128\}$, with a twofold aim: to check that all three reproduce the equilibrium physics and to compare their dynamical behaviour. At equilibrium the three agree, as they must when sampling the same Boltzmann distribution. The estimate of $T_c$ obtained from the Binder cumulant crossings is compatible with Onsager's value (with $U_4^*\approx 0{,}61$), and the FSS collapse of $\chi$ and $\xi_L/L$ through the search for the critical exponents confirms that the data belong to the universality class of the two-dimensional Ising model. The only caveat appears near $T_c$, in the observables most sensitive to autocorrelation, where Metropolis and Glauber feel the critical slowing down, so that the clean convergence of $T_c$ with $L$ is seen only with Wolff and the peak of the collapsed susceptibility shifts to $x\approx 1$ for the local algorithms, against $x\approx 1{,}5$ for Wolff, which is the value we take as reliable.

The underlying difference between the algorithms lies in the dynamics, measured through the integrated autocorrelation time $\tau_{\text{int}}$. Metropolis and Glauber suffer critical slowing down, with a peak of $\tau_{\text{int}}$ that grows with $L$ around $T_c$ consistently with $z\approx 2$; within this class, Glauber is somewhat slower than Metropolis (by a factor between $1{,}5$ and $2$), because Metropolis accepts the flips with a probability greater than or equal to that of Glauber for every $\Delta E$ and decorrelates sooner, which is a general result. Wolff, by contrast, practically eliminates the critical slowing down: its $\tau_{\text{int}}$ is of unit order over the whole range ($z\approx 0{,}017$), so that near $T_c$ obtaining an independent sample costs several orders of magnitude less than with the local algorithms, an advantage that would become decisive for larger lattices.

As a natural continuation, it would be interesting to extend the comparison to the Ising model in three dimensions, where there is no exact solution and the exponents are known only numerically, now that we know that the program together with the data analysis works satisfactorily, as shown by comparing the results obtained against the exact results of the two-dimensional Ising model.
\end{otherlanguage}